\chapter{Actions à deux catégories, sociétés holding, fonds fermés et SPAC}\label{ch:s1:share-classes-holding-companies-closed-end-funds-spacs}

Un fonds fermé détient 100~millions de dollars d'actions cotées et se négocie à 85~millions de dollars. Ses parts ne peuvent pas être rachetées contre les actions, si bien que rien
n'oblige les deux valeurs à se rejoindre~: la décote peut durer des années, se resserrer quand le sentiment se retourne, ou s'élargir. Acheter le fonds et vendre ses titres
ne verrouille quinze cents par dollar que si quelque chose referme l'écart. Sur le modèle synthétique de ce chapitre, la transaction rapporte 0.87~\% par an si la décote tend
à revenir vers sa moyenne, avec une détention médiane de trois ans et une probabilité de deux sur cinq de perdre~; avec des catalyseurs (une \omterm{def:s1:merger-arbitrage:arb}{offre publique} de rachat, une liquidation, un
investisseur activiste) survenant à raison de 0.3 par an, elle rapporte 5.0~\% par an et se clôt au bout de 1.4 an en médiane. La valeur relative structurelle désigne la famille des transactions dans lesquelles un
contrat ou une structure lie deux prix entre eux~: \omterm{def:s1:share-classes-holding-companies-closed-end-funds-spacs:dual}{actions à deux catégories}, sociétés holding, fonds fermés et SPAC. La
construction correspond à \texttt{firm.structrv}.

\section{Actions à deux catégories}

\begin{definition}[Action à deux catégories]\label{def:s1:share-classes-holding-companies-closed-end-funds-spacs:dual}
Une \emph{action à deux catégories}\index{action à deux catégories} est une structure qui donne à une société deux catégories ou plus d'actions ordinaires cotées, ayant le même droit sur ses flux de trésorerie mais des droits différents
(en général les droits de vote), de sorte que les catégories devraient se négocier à des niveaux proches les unes des autres et que leur spread traduise la valeur de la différence de droits, de liquidité
et d'éligibilité aux indices.
\end{definition}

Les \omterm{def:s1:share-classes-holding-companies-closed-end-funds-spacs:dual}{actions à deux catégories} forment la paire structurelle la plus pure~: la même société, les mêmes dividendes, un spread qui devrait être faible et stationnaire. Elles sont la
contrepartie réelle des jumeaux disséminés du chapitre~4, dont le spread de cours logarithmique revenait vers zéro avec une constante de temps de dix jours et un écart type de 2~\%, et que la
\omterm{def:s1:pairs-and-baskets:distance}{méthode des distances} a repérés et négociés avec un ratio de Sharpe de 5 après coûts. Les spreads réels sont déterminés par ce qui différencie les catégories~: les droits de vote
comptent dans les batailles pour le contrôle, l'une des catégories peut figurer dans un indice et l'autre non (chapitre~10), et l'une peut être beaucoup moins liquide. Un livre d'\omterm{def:s1:share-classes-holding-companies-closed-end-funds-spacs:dual}{actions à deux catégories} négocie le
spread par rapport à son propre historique et clôt les positions avant les événements susceptibles de modifier la valeur du droit de vote.

\section{Décotes des sociétés holding}

\begin{definition}[Décote de holding]\label{def:s1:share-classes-holding-companies-closed-end-funds-spacs:holding}
Une \emph{décote de holding}\index{décote de holding} est l'écart dont la valeur de marché d'une société holding est inférieure à la valeur de ses participations dans d'autres sociétés cotées et
non cotées (après déduction de sa dette), exprimé en fraction de cette valeur.
\end{definition}

Une société holding qui détient des participations cotées peut être couverte à l'identique~: acheter la holding et vendre ses participations au prorata. Ce qui reste est la décote,
qui a des raisons d'exister (impôts sur la cession des participations, frais de la holding, gouvernance) et peu de forces pour la résorber~; à la différence d'un fonds fermé, une holding
se liquide rarement. La société holding du chapitre se négocie à 70~cents par dollar, avec une décote qui revient vers 80~cents selon une demi-vie de trois
ans et sans catalyseur~: la transaction couverte rapporte 0.94~\% par an, se clôt dans les trois ans sur seulement 2~\% des trajectoires et perd sur 40~\% d'entre elles.

\section{Décotes des fonds fermés}

\begin{definition}[Décote de fonds fermé]\label{def:s1:share-classes-holding-companies-closed-end-funds-spacs:cef}
Une \emph{décote de fonds fermé}\index{décote de fonds fermé} est l'écart dont le cours d'un fonds fermé est inférieur à sa valeur liquidative par part, exprimé en fraction de la
valeur liquidative~; un fonds fermé émet un nombre fixe de parts qui se négocient en bourse et ne peuvent pas être rachetées à la valeur liquidative.
\end{definition}

L'énigme du fonds fermé comporte deux volets. Les fonds se négocient à une décote la plupart du temps, et ces décotes varient. Lee, Shleifer et Thaler ont attribué le
mouvement au sentiment des investisseurs particuliers~: les décotes de fonds différents évoluent ensemble, de nouveaux fonds se créent lorsque les fonds établis se vendent près de leur valeur d'actif ou au-dessus,
et les décotes se resserrent quand les petites capitalisations progressent. Pontiff a constaté que les écarts sont plus importants là où l'arbitrage coûte plus cher~: portefeuilles difficiles à répliquer, dividendes
faibles, petits fonds et taux d'intérêt élevés, qui expliquent ensemble un quart de la dispersion transversale des mauvaises évaluations.

\begin{center}
\small
\setlength{\tabcolsep}{4pt}%
\begin{tabular}{@{}p{4.6cm}cccc@{}}
\toprule
& \multicolumn{3}{c}{fonds fermé à 85 cents} & société holding\\
catalyseurs par an & aucun & 0.1 & 0.3 & à 70 cents, aucun\\
\midrule
rendement moyen par transaction~; par an & 1.9~\%~; 0.87~\% & 4.4~\%~; 2.23~\% & 7.9~\%~; 5.04~\% & 2.8~\%~; 0.94~\%\\
durée médiane de détention (années) & 2.96 & 2.13 & 1.37 & 3.00\\
clôturées dans les trois ans & 50~\% & 63~\% & 79~\% & 2~\%\\
clôturées par un catalyseur & 0~\% & 19~\% & 46~\% & 0~\%\\
5\up{e} centile du rendement~; part des pertes & $-14.7~\%$~; 40~\% & $-13.3~\%$~; 30~\% & $-10.4~\%$~; 17~\% & $-13.9~\%$~; 40~\%\\
\bottomrule
\end{tabular}
\end{center}

Le modèle (\cref{lst:s1:share-classes-holding-companies-closed-end-funds-spacs:discount}) est une décote à retour à la moyenne, exprimée en logarithmes~: le fonds part à 85 cents
par dollar ($-16.3~\%$) et revient vers 90 cents ($-10.5~\%$) avec une demi-vie de deux ans et une volatilité de 8~\% par an. La transaction achète le fonds et vend
à découvert son portefeuille, supporte un carry de 1.3~\% par an (les frais du fonds et l'emprunt de la couverture), et se clôt lorsque la décote atteint 5~\% ou au bout de trois ans.
Sans catalyseur, l'essentiel du resserrement attendu va vers la décote moyenne, et non vers zéro, et le carry en absorbe une grande part~: la transaction est un gain lent et incertain de
2~\% sur trois ans. Un catalyseur change le calcul~: la décote se réduit à 2~\% d'un coup, le capital est libéré plus tôt, et le rendement annuel
se multiplie par près de six avec 0.3 catalyseur par an (\cref{fig:s1:share-classes-holding-companies-closed-end-funds-spacs:exit}). C'est pourquoi les transactions sur décote sont si souvent
des transactions activistes~: l'investisseur capable d'imposer une \omterm{def:s1:merger-arbitrage:arb}{offre publique} de rachat ou une liquidation détient le catalyseur.

\begin{omfigure}
\begin{tikzpicture}
\begin{axis}[ybar,bar width=7pt,width=11cm,height=5.0cm,xlabel={\small années jusqu'à la clôture de la transaction},ylabel={\small part des trajectoires (\%)},
  tick label style={font=\footnotesize},xmin=0,xmax=3.1,ymin=0,ymax=60,grid=major,grid style={gray!25},
  legend style={font=\scriptsize,at={(0.03,0.97)},anchor=north west},table/col sep=comma]
\addplot[fill=omThm!60,draw=omThm] table[x=years,y=no_catalyst]{figdata/strategies-1/12-share-classes-holding-companies-closed-end-funds-spacs/exit.csv};
\addplot[fill=omDef!60,draw=omDef] table[x=years,y=catalyst]{figdata/strategies-1/12-share-classes-holding-companies-closed-end-funds-spacs/exit.csv};
\legend{sans catalyseur,0.3 catalyseurs par an}
\end{axis}
\end{tikzpicture}

\omcaption{La transaction sur \omterm{def:s1:share-classes-holding-companies-closed-end-funds-spacs:cef}{décote de fonds fermé}~: moment de sa clôture (la décote atteint 5~\%, un catalyseur survient ou trois ans s'écoulent), sur 10\,000
trajectoires simulées~; la dernière barre regroupe les transactions encore ouvertes à trois ans. Données~: \texttt{s1\_structrv.fund}.}
\label{fig:s1:share-classes-holding-companies-closed-end-funds-spacs:exit}
\end{omfigure}

\section{Sociétés à chèque en blanc et leur valeur de compte séquestre}

\begin{definition}[Special purpose acquisition company, valeur du compte séquestre]\label{def:s1:share-classes-holding-companies-closed-end-funds-spacs:spac}
Un \emph{special purpose acquisition company}\index{special purpose acquisition company} (SPAC) est une société-écran cotée qui lève des fonds lors d'une introduction en bourse initiale, les place sur un compte séquestre et dispose
d'un délai limité pour fusionner avec une société non cotée~; ses actionnaires peuvent racheter leurs actions contre leur part du compte séquestre plutôt que de participer à la fusion.
La \emph{valeur du compte séquestre}\index{valeur du compte séquestre} (trust value) par action est ce montant de rachat, soit les fonds levés augmentés des intérêts accumulés par le compte.
\end{definition}

Le droit de rachat fait d'une action de SPAC un plancher plus une option~: lors du vote, le détenteur peut prendre la \omterm{def:s1:share-classes-holding-companies-closed-end-funds-spacs:spac}{valeur du compte séquestre} ou conserver une part de la société issue de la fusion. Une
action achetée en dessous de la \omterm{def:s1:share-classes-holding-companies-closed-end-funds-spacs:spac}{valeur du compte séquestre} a un rendement minimal garanti (si le compte est sûr) et un call gratuit sur la fusion. Achetée à 9.80~\$ avec un compte séquestre de
10.00~\$ rapportant 4~\% par an et le vote étant à un an, le plancher s'élève à 10.41~\$, soit un rendement de 6.2~\% (\cref{lst:s1:share-classes-holding-companies-closed-end-funds-spacs:spac}).
La valeur du call dépend de la société issue de la fusion. Klausner, Ohlrogge et Ruan ont constaté que les SPAC qui levaient 10.00~\$ par action conservaient bien moins de trésorerie nette par action
lors de leur fusion (5.70~\$ en médiane pour celles qui fusionnaient de janvier 2019 à juin 2020), et que les actionnaires qui gardaient leurs titres jusqu'à la fusion subissaient de fortes
pertes tandis que les promoteurs en tiraient profit. L'hypothèse du chapitre suit ce constat~: une valeur de fusion lognormale autour de 7~\$ avec une volatilité de 60~\%. L'arbitragiste
qui rachète lorsque la valeur de fusion est inférieure au compte séquestre (sur 74~\% des trajectoires) gagne 20.2~\% en moyenne, soit le plancher plus un call sur une issue volatile~; un détenteur
qui ne rachète jamais perd 13.9~\% en moyenne et 27.8~\% à la médiane.

\begin{dated}{2026-09}{Règles de la SEC pour les SPAC}\label{dat:s1:share-classes-holding-companies-closed-end-funds-spacs:sec}
Le 24 janvier 2024, la SEC a adopté des règles et des modifications applicables aux introductions en bourse initiales de SPAC et aux opérations de fusion après introduction en bourse (de-SPAC, soit les fusions), imposant entre autres
une information renforcée sur les conflits d'intérêts, la rémunération des promoteurs et la dilution, davantage d'informations sur la société cible, et encadrant
l'usage des projections, afin d'aligner les SPAC de plus près sur les introductions en bourse traditionnelles.
\end{dated}

\section{Valeur relative structurelle, en général}

Les quatre stratégies partagent une même forme. Une structure lie un cours à une valeur qui peut être observée ou répliquée (l'autre catégorie, les participations, la valeur liquidative,
le compte séquestre). L'écart existe parce que personne ne peut le refermer à peu de frais~: votes, impôts, impossibilité de racheter, délai avant un vote. La stratégie tire profit de la
réduction de l'écart et perçoit du carry pendant qu'elle attend. Ses risques sont l'élargissement de l'écart (le sentiment, comme l'ont montré Lee, Shleifer et Thaler pour les fonds), l'échec de la couverture (un
portefeuille qu'il est impossible de vendre à découvert à l'identique, comme l'a montré Pontiff), et le temps. La question pratique est toujours la même~: qu'est-ce qui referme l'écart, et quand~? Avec un
catalyseur, la stratégie est un événement~; sans lui, c'est un carry trade lent et incertain.

\section{Fiches de stratégie}

\begin{strategyfile}{Spread entre catégories d'actions}\label{strat:s1:share-classes-holding-companies-closed-end-funds-spacs:class}
\sfield{Qui paie, et pourquoi} Des traders qui ont besoin de la liquidité d'une catégorie ou de son appartenance à un indice, et qui écartent son cours de celui de l'autre.
\sfield{Instruments et plateformes} Les deux catégories cotées d'une même société.
\sfield{Signal} Le spread logarithmique comparé à son historique propre, net de la valeur de la différence de droits de vote.
\sfield{Taille et exécution} Position longue sur la catégorie bon marché, courte sur la chère~; petites transactions fréquentes.
\sfield{Coûts} Coût de l'emprunt et fourchette sur la catégorie la moins liquide.
\sfield{Comment elle s'éteint} Les batailles pour le contrôle qui réévaluent les droits de vote~; les changements d'indice~; la conversion ou l'effondrement de la structure.
\sfield{Horizon, capacité, infrastructure} De quelques jours à quelques semaines~; capacité fixée par la catégorie la moins liquide.
\sfield{Backtest honnête} Les cours des deux catégories au même instant~; les événements sur le capital présents dans l'échantillon.
\sfield{Sources} La simulation jumelle du chapitre~4~; aucun historique de performance public vérifié.
\end{strategyfile}

\begin{strategyfile}{Stratégie de décote de holding}\label{strat:s1:share-classes-holding-companies-closed-end-funds-spacs:holding}
\sfield{Qui paie, et pourquoi} Des investisseurs qui fuient la structure, la fiscalité ou la gouvernance de la holding, et qui la vendent en dessous de la somme de ses participations.
\sfield{Instruments et plateformes} La holding, en position longue~; ses participations cotées, en position courte.
\sfield{Signal} La décote comparée à son historique et à ses causes (fiscalité, coûts, changements de gouvernance).
\sfield{Taille et exécution} Couverte par les participations cotées~; la partie non cotée est laissée en risque ou estimée.
\sfield{Coûts} Coût de l'emprunt sur les participations~; carry pendant plusieurs années.
\sfield{Comment elle s'éteint} Décotes qui persistent ou s'élargissent~; pas de catalyseur.
\sfield{Horizon, capacité, infrastructure} Plusieurs années~; le portefeuille publié de la holding.
\sfield{Backtest honnête} Participations telles que publiées à chaque date~; la partie non cotée évaluée de façon prudente.
\sfield{Sources} La simulation de ce chapitre (0.94~\% par an sans catalyseur)~; aucun historique de performance public vérifié.
\end{strategyfile}

\begin{strategyfile}{Retour à la moyenne de la décote des fonds fermés}\label{strat:s1:share-classes-holding-companies-closed-end-funds-spacs:cef}
\sfield{Qui paie, et pourquoi} Des investisseurs particuliers dont le sentiment fait varier les décotes, et qui vendent des fonds lorsque les décotes sont élevées.
\sfield{Instruments et plateformes} Fonds fermés, en position longue~; leurs portefeuilles ou un proxy, en position courte.
\sfield{Signal} La décote comparée à sa moyenne de long terme et aux décotes des autres fonds.
\sfield{Taille et exécution} De nombreux fonds, chacun couvert~; sortie proche de la décote moyenne.
\sfield{Coûts} Les frais du fonds pendant la détention~; les coûts de couverture lorsque le portefeuille est difficile à répliquer.
\sfield{Comment elle s'éteint} Un sentiment qui élargit ensemble toutes les décotes.
\sfield{Horizon, capacité, infrastructure} De quelques mois à plusieurs années~; valeurs liquidatives quotidiennes.
\sfield{Backtest honnête} Valeurs liquidatives telles que publiées (avec leur décalage)~; frais et coûts de couverture.
\sfield{Sources} Lee, Shleifer et Thaler (1991)~; Pontiff (1996).
\end{strategyfile}

\begin{strategyfile}{Stratégie de décote à catalyseur activiste}\label{strat:s1:share-classes-holding-companies-closed-end-funds-spacs:activist}
\sfield{Qui paie, et pourquoi} Gérants de fonds et conseils d'administration qui laissent une décote ouverte jusqu'à ce qu'ils soient contraints de la refermer.
\sfield{Instruments et plateformes} Fonds fermés et sociétés de holding aux décotes élevées et à la gouvernance contestable.
\sfield{Signal} La décote, et la faisabilité d'une \omterm{def:s1:merger-arbitrage:arb}{offre publique} imposée, d'une liquidation ou d'une conversion.
\sfield{Taille et exécution} Positions assez importantes pour peser sur l'issue~; campagnes de procuration.
\sfield{Coûts} Coûts juridiques et de campagne~; longues durées de détention si la campagne échoue.
\sfield{Comment elle s'éteint} Mesures défensives~; autres actionnaires qui ne soutiennent pas la campagne.
\sfield{Horizon, capacité, infrastructure} De un à trois ans~; capacité juridique et de gouvernance.
\sfield{Backtest honnête} Les campagnes qui ont échoué comme celles qui ont réussi.
\sfield{Sources} La simulation de ce chapitre (5.0~\% par an avec 0.3 catalyseur par an, contre 0.87~\% sans)~; aucun historique de performance public vérifié.
\end{strategyfile}

\begin{strategyfile}{Plancher de la valeur du compte séquestre des SPAC}\label{strat:s1:share-classes-holding-companies-closed-end-funds-spacs:spac}
\sfield{Qui paie, et pourquoi} Des investisseurs qui vendent des actions de SPAC en dessous de leur \omterm{def:s1:share-classes-holding-companies-closed-end-funds-spacs:spac}{valeur du compte séquestre}, et la valeur optionnelle de la fusion.
\sfield{Instruments et plateformes} Actions de SPAC (et parfois warrants) avant le vote sur la fusion.
\sfield{Signal} Cours inférieur à la \omterm{def:s1:share-classes-holding-companies-closed-end-funds-spacs:spac}{valeur du compte séquestre} majorée du coupon couru~; temps restant avant le vote ou l'échéance.
\sfield{Taille et exécution} Achat sous la \omterm{def:s1:share-classes-holding-companies-closed-end-funds-spacs:spac}{valeur du compte séquestre}, rachat au vote sauf si l'action issue de la fusion vaut davantage~; de nombreux SPAC.
\sfield{Coûts} Faibles~; capital immobilisé jusqu'au vote.
\sfield{Comment elle s'éteint} Comptes séquestres qui rapportent moins que prévu~; prorogations qui allongent l'attente~; règles et marchés qui modifient le produit.
\sfield{Horizon, capacité, infrastructure} De quelques mois à deux ans~; conditions du compte séquestre et du rachat tirées des dossiers déposés.
\sfield{Backtest honnête} Rachat à la \omterm{def:s1:share-classes-holding-companies-closed-end-funds-spacs:spac}{valeur du compte séquestre}, et non au cours après fusion~; l'ensemble de l'échantillon de SPAC.
\sfield{Sources} Klausner, Ohlrogge et Ruan (2022) sur la trésorerie nette et les pertes après fusion~; la simulation de ce chapitre (plancher de 6.2~\%).
\end{strategyfile}

\section{Tutoriel~: quatre-vingt-cinq cents pour un dollar}\label{tut:s1:share-classes-holding-companies-closed-end-funds-spacs:tut}

\begin{tutorial}
\textbf{Objectif.} Simuler une décote qui se résorbe lentement et qui peut être refermée par un catalyseur, la négocier, et évaluer un SPAC comme un plancher et un call. \textbf{État final~:}
le tableau et la \cref{fig:s1:share-classes-holding-companies-closed-end-funds-spacs:exit}.
\begin{tutsteps}
\item \textbf{La décote et la stratégie}~: un AR(1) en logarithme avec des catalyseurs de Poisson~; sortie à un niveau, sur un catalyseur ou à l'horizon.
\omcode{code/firm/structrv/firm_structrv.py}{27}{52}{Les trajectoires de la décote et la stratégie couverte.}\label{lst:s1:share-classes-holding-companies-closed-end-funds-spacs:discount}
\item \textbf{Le SPAC}~: le plancher du compte séquestre et l'option sur la fusion.
\omcode{code/firm/structrv/firm_structrv.py}{55}{60}{Le gain d'une action de SPAC au moment du vote.}\label{lst:s1:share-classes-holding-companies-closed-end-funds-spacs:spac}
\item \textbf{Exécuter} \texttt{fund(0.0)}, \texttt{fund(0.1)}, \texttt{fund(0.3)}, \texttt{holding()}, \texttt{spac()} et \texttt{fig\_structrv.py}.
\end{tutsteps}
\textbf{Pour aller plus loin.} Faire varier ensemble les décotes de plusieurs fonds (le sentiment commun de Lee, Shleifer et Thaler) et dimensionner un portefeuille de ces fonds~; ajouter
des warrants au SPAC~; faire dépendre l'arrivée du catalyseur de la largeur de la décote.
\end{tutorial}

\section{Construction~: valeur relative structurelle}\label{bld:s1:share-classes-holding-companies-closed-end-funds-spacs:structrv}

\begin{build}
\bfield{Rôle} Dynamique de la décote avec catalyseurs, la stratégie de décote couverte, et le gain du SPAC, plancher plus option.
\bfield{Interface} \texttt{simulate\_discount(d0, mean, half\_life, vol, days, catalyst\_rate, catalyst\_level, paths, rng)},
\texttt{discount\_trade(paths, catalyst\_day, close\_level, carry, horizon)}, \texttt{spac\_payoff(trust\_at\_vote, merged\_value)},
\texttt{spac\_return(price, trust, rate, days, merged\_value)}.
\bfield{Règles} Décotes en logarithmes~; carry imputé pour la durée de détention~; sorties au premier de ces événements~: niveau, catalyseur ou horizon.
\bfield{Tests d'acceptation} \texttt{code/firm/structrv/tests/}~: retour déterministe à la moyenne avec volatilité nulle~; catalyseurs qui referment la décote~; la sortie de la
stratégie et son rendement, calculés à la main~; le gain du SPAC, calculé à la main.
\bfield{Pour aller plus loin} Décotes corrélées entre fonds~; campagnes activistes avec coûts~; warrants du SPAC et votes de prorogation.
\end{build}

\begin{omsources}
\item C.~M.~C.~Lee, A.~Shleifer et R.~H.~Thaler, ``Investor sentiment and the closed-end fund puzzle'', \emph{Journal of Finance} 46(1), 1991.
\item J.~Pontiff, ``Costly arbitrage: evidence from closed-end funds'', \emph{Quarterly Journal of Economics} 111(4), 1996.
\item M.~Klausner, M.~Ohlrogge et E.~Ruan, ``A sober look at SPACs'', \emph{Yale Journal on Regulation} 39, 2022.
\item US SEC, communiqué de presse 2024-8, SPAC rules, janvier 2024.
\end{omsources}

\section{Exercices}

\begin{exercise}[$\star$]\label{exo:s1:share-classes-holding-companies-closed-end-funds-spacs:1}
Exprimer des décotes de 85, 90 et 95 cents pour un dollar en logarithmes.
\end{exercise}

\begin{exercise}[$\star$]\label{exo:s1:share-classes-holding-companies-closed-end-funds-spacs:2}
Une transaction achetée à 85 cents clôture à 95 cents après un an, avec un carry de 1.3~\%. Quel est son rendement logarithmique~?
\end{exercise}

\begin{exercise}[$\star$]\label{exo:s1:share-classes-holding-companies-closed-end-funds-spacs:3}
Un SPAC se négocie à 9.80~\$ avec un compte séquestre de 10.00~\$ qui rapporte 4~\% par an, et un vote dans un an. Quelle est la \omterm{def:s1:share-classes-holding-companies-closed-end-funds-spacs:spac}{valeur du compte séquestre} (trust value) au moment du vote, et quel est le rendement plancher~?
\end{exercise}

\begin{exercise}[$\star\star$]\label{exo:s1:share-classes-holding-companies-closed-end-funds-spacs:4}
Avec une demi-vie de deux ans, quelle fraction de l'écart entre la décote actuelle et sa moyenne subsiste au bout d'un an~? Pourquoi la stratégie sans catalyseur est-elle si
lente~?
\end{exercise}

\begin{exercise}[$\star\star$]\label{exo:s1:share-classes-holding-companies-closed-end-funds-spacs:5}
Pourquoi une \omterm{def:s1:share-classes-holding-companies-closed-end-funds-spacs:cef}{décote des fonds fermés} est-elle plus difficile à arbitrer qu'un écart entre deux catégories d'actions~?
\end{exercise}

\begin{exercise}[$\star\star$]\label{exo:s1:share-classes-holding-companies-closed-end-funds-spacs:6}
Expliquer pourquoi un SPAC acheté sous sa \omterm{def:s1:share-classes-holding-companies-closed-end-funds-spacs:spac}{valeur du compte séquestre} (trust value) équivaut à un plancher plus une option d'achat, et ce qui peut faire échouer ce plancher.
\end{exercise}

\begin{exercise}[$\star\star\star$]\label{exo:s1:share-classes-holding-companies-closed-end-funds-spacs:7}
\emph{Programmation.} Exécuter \texttt{fund(0.1)} et comparer le rendement annuel et la durée de détention avec \texttt{fund(0.0)} et \texttt{fund(0.3)}. Déterminer combien il vaudrait la
peine de payer par an pour un catalyseur qui arrive à 0.3 par an~?
\end{exercise}

\begin{exercise}[$\star\star\star$]\label{exo:s1:share-classes-holding-companies-closed-end-funds-spacs:8}
\emph{Trouver la faille.} «~Le fonds se négocie avec une décote de 15~\% sur ses actifs~; l'achat du fonds combiné à la vente à découvert de son portefeuille garantit un gain de 15~\%.~»
\end{exercise}

\section{Problème~: Quatre-vingt-cinq cents pour un dollar}

\begin{problem}[{Problème du week-end --- ce qui referme l'écart}]\label{pb:s1:share-classes-holding-companies-closed-end-funds-spacs:1}
Les décotes et le SPAC simulés du chapitre, ainsi que les données publiques.

\medskip\noindent\textbf{Partie I --- Structures.}
\begin{enumerate}
\item Définir une \omterm{def:s1:share-classes-holding-companies-closed-end-funds-spacs:dual}{action à deux catégories}, une \omterm{def:s1:share-classes-holding-companies-closed-end-funds-spacs:holding}{décote de holding} et une \omterm{def:s1:share-classes-holding-companies-closed-end-funds-spacs:cef}{décote des fonds fermés}.
\item Pourquoi ces écarts existent-ils, et pourquoi persistent-ils~?
\item Mettre en relation les catégories d'actions avec les jumeaux du chapitre~4.
\item Qu'ont établi Lee, Shleifer et Thaler, ainsi que Pontiff~?
\end{enumerate}

\medskip\noindent\textbf{Partie II --- La stratégie sur décote.}
\begin{enumerate}[resume]
\item Décrire le modèle de la décote du fonds et la stratégie.
\item Donner les résultats sans catalyseur.
\item Donner les résultats avec des catalyseurs à 0.1 et 0.3 par an.
\item Que montre le résultat de la holding~?
\end{enumerate}

\medskip\noindent\textbf{Partie III --- SPAC.}
\begin{enumerate}[resume]
\item Définir un SPAC et sa \omterm{def:s1:share-classes-holding-companies-closed-end-funds-spacs:spac}{valeur du compte séquestre} (trust value).
\item Calculer le plancher du SPAC du chapitre.
\item Qu'ont établi Klausner, Ohlrogge et Ruan~?
\item Que dit la simulation sur le rachat et la détention~?
\end{enumerate}

\medskip\noindent\textbf{Partie IV --- Le verdict.}
\begin{enumerate}[resume]
\item Énoncer le \emph{résultat nommé}~: le rendement espéré et le temps de convergence d'une stratégie sur décote, avec et sans catalyseur.
\item Qu'ont en commun les quatre structures~?
\item Pourquoi les stratégies sur décote sont-elles si souvent des stratégies d'activiste~?
\item Qu'ont changé pour les SPAC les règles de la SEC de 2024~?
\item Comment couvrir un fonds dont le portefeuille ne peut pas être vendu à découvert~?
\item Quelle fiche de stratégie lancer avec un capital patient~?
\item Quel est le risque principal d'un portefeuille de stratégies sur décote~?
\item En une phrase~: que procure l'arbitrage de valeur relative structurelle~?
\end{enumerate}
\end{problem}

\section{Questions d'entretien}

\begin{interviewq}[$\star$ \iqroles{chercheur, trader}]\label{iq:s1:share-classes-holding-companies-closed-end-funds-spacs:1}
Pourquoi un fonds fermé peut-il se négocier sous sa valeur liquidative pendant des années~?
\end{interviewq}

\begin{interviewq}[$\star\star$ \iqroles{trader}]\label{iq:s1:share-classes-holding-companies-closed-end-funds-spacs:2}
Comment négocieriez-vous deux catégories d'actions d'une même société~?
\end{interviewq}

\begin{interviewq}[$\star\star$ \iqroles{chercheur}]\label{iq:s1:share-classes-holding-companies-closed-end-funds-spacs:3}
Comment évalueriez-vous une action de SPAC avant le vote de fusion~?
\end{interviewq}

\begin{interviewq}[$\star\star$ \iqroles{risque}]\label{iq:s1:share-classes-holding-companies-closed-end-funds-spacs:4}
Quels sont les risques d'un portefeuille de positions couvertes sur fonds fermés~?
\end{interviewq}

\begin{interviewq}[$\star\star$ \iqroles{trader, chercheur}]\label{iq:s1:share-classes-holding-companies-closed-end-funds-spacs:5}
Une société de holding se négocie à 70~\% de la valeur de ses parts. Qu'est-ce qui vous amènerait à l'acheter~?
\end{interviewq}

\begin{interviewq}[$\star\star\star$ \iqroles{chercheur}]\label{iq:s1:share-classes-holding-companies-closed-end-funds-spacs:6}
Une décote suit un processus d'Ornstein--Uhlenbeck vers $\bar d$ avec une vitesse $\kappa$, et un catalyseur la ramène à $d_c$ au taux de Poisson $\lambda$. Écrire la
décote espérée à l'instant $t$ et le gain espéré d'une transaction conservée jusqu'à $t$.
\end{interviewq}
